\documentclass[reqno]{amsart}
\usepackage[T1]{fontenc}
\usepackage{amsmath,amssymb,amsthm}
\usepackage{xr-hyper}
\usepackage[hidelinks]{hyperref}
\externaldocument[sd-]{20261125T000000--slope-method__workbook}
\externaldocument[cg-]{20261125T014100--congruence-modules__workbook}
\externaldocument[da-]{20261126T030300--the-data__workbook}
\externaldocument[ob-]{20261126T040100--spaces-forms-and-twists__workbook}
\externaldocument[dl-]{20261126T030500--divisibility-of-the-composite-lift__workbook}
\externaldocument[tw-]{20261126T031000--twists-with-few-congruences__workbook}
\externaldocument[col-]{20261126T030800--the-cost-of-one-lift__workbook}
\externaldocument[l1-]{20261126T050800--compressions-and-false-lifts__workbook}
\externaldocument[l3-]{20261126T050900--hecke-values-as-norms-of-lifts__workbook}
\externaldocument[l4-]{20261126T051000--the-bound-is-attained-at-many-places__workbook}
\externaldocument[c3-]{20261126T050600--the-exchange-identity__workbook}

\newcounter{chap}
% Chapter numbers, generated from 20261125T012000--grh-workbook__guide.org
% by scripts/make-chapters.py.  Do not edit: change the order of the
% index and rerun the script (./preview-site.sh and CI do this).
%
% Usage in a sheet, after \newcounter{chap}:
%   % Chapter numbers, generated from 20261125T012000--grh-workbook__guide.org
% by scripts/make-chapters.py.  Do not edit: change the order of the
% index and rerun the script (./preview-site.sh and CI do this).
%
% Usage in a sheet, after \newcounter{chap}:
%   % Chapter numbers, generated from 20261125T012000--grh-workbook__guide.org
% by scripts/make-chapters.py.  Do not edit: change the order of the
% index and rerun the script (./preview-site.sh and CI do this).
%
% Usage in a sheet, after \newcounter{chap}:
%   \input{chapters}\chapnumber{<sheet ID>}
\makeatletter
\providecommand\chapters@set[2]{\@namedef{chapters@#1}{#2}}
\chapters@set{20261125T001000}{1}% zeros-to-ramp-sums
\chapters@set{20261125T001100}{2}% inert-primes-and-estimate-b
\chapters@set{20261125T005000}{3}% galois-fourier-and-averaging
\chapters@set{20261125T005500}{4}% columns-from-a-sphere
\chapters@set{20261125T000000}{5}% slope-method
\chapters@set{20261126T030100}{6}% a-twist-for-each-prime
\chapters@set{20261125T003000}{7}% the-chain-and-the-square-root
\chapters@set{20261125T010000}{8}% rank-lemma-and-zero-estimate
\chapters@set{20261125T010500}{9}% poincare-series-and-near-orthogonality
\chapters@set{20261124T223000}{10}% elliptic-curves-for-the-proof
\chapters@set{20261125T009500}{11}% fock-model-and-theta-intertwiners
\chapters@set{20261124T223100}{12}% complex-multiplication-and-inert-reduction
\chapters@set{20261125T008000}{13}% one-cm-determinant
\chapters@set{20261124T214500}{14}% supersingular-theta-divisibility
\chapters@set{20261124T214600}{15}% the-isotropic-plane
\chapters@set{20261125T013700}{16}% pairings-on-a-formal-group-of-height-two
\chapters@set{20261125T013500}{17}% hecke-nondegeneracy-on-the-tree
\chapters@set{20261126T030200}{18}% the-bound-is-attained
\chapters@set{20261125T009000}{19}% kronecker-limit-and-chowla-selberg
\chapters@set{20261125T002000}{20}% definite-theta-lifts
\chapters@set{20261125T015500}{21}% siegel-weil-in-the-smallest-case
\chapters@set{20261125T002100}{22}% definite-theta-lifts-and-norms
\chapters@set{20261125T016500}{23}% hecke-operators-and-theta-lifts
\chapters@set{20261125T016700}{24}% divisibility-along-the-chain
\chapters@set{20261125T004000}{25}% denominators-and-sha
\chapters@set{20261125T004100}{26}% tate-twists-and-sha
\chapters@set{20261126T020430}{27}% the-minus-part-by-gauss-sums
\chapters@set{20261125T014100}{28}% congruence-modules
\chapters@set{20261125T014500}{29}% eisenstein-congruences-and-selmer-classes
\chapters@set{20261125T014700}{30}% pseudocharacters-and-the-length-bound
\chapters@set{20261125T011000}{31}% galois-representations-black-box
\chapters@set{20261126T030300}{32}% the-data
\chapters@set{20261126T020210}{33}% bessel-poincare-series-and-their-gram-operators
\chapters@set{20261126T030400}{34}% the-seesaw-identity
\chapters@set{20261126T030500}{35}% divisibility-of-the-composite-lift
\chapters@set{20261126T030600}{36}% theorem-a-short-forms
\chapters@set{20261126T030700}{37}% the-determinant-of-a-chain
\chapters@set{20261126T030800}{38}% the-cost-of-one-lift
\chapters@set{20261126T030900}{39}% denominators-of-projections-to-lifts
\chapters@set{20261126T031000}{40}% twists-with-few-congruences
\chapters@set{20261126T031100}{41}% theorem-c-and-the-main-inequality
\providecommand\chapnumber[1]{%
  \@ifundefined{chapters@#1}%
    {\PackageError{chapters}{Sheet #1 is not listed in the index}%
      {List it in 20261125T012000--grh-workbook__guide.org\MessageBreak
       and run scripts/make-chapters.py.}}%
    {\edef\thechap{\@nameuse{chapters@#1}}}}
\makeatother
\chapnumber{<sheet ID>}
\makeatletter
\providecommand\chapters@set[2]{\@namedef{chapters@#1}{#2}}
\chapters@set{20261125T001000}{1}% zeros-to-ramp-sums
\chapters@set{20261125T001100}{2}% inert-primes-and-estimate-b
\chapters@set{20261125T005000}{3}% galois-fourier-and-averaging
\chapters@set{20261125T005500}{4}% columns-from-a-sphere
\chapters@set{20261125T000000}{5}% slope-method
\chapters@set{20261126T030100}{6}% a-twist-for-each-prime
\chapters@set{20261125T003000}{7}% the-chain-and-the-square-root
\chapters@set{20261125T010000}{8}% rank-lemma-and-zero-estimate
\chapters@set{20261125T010500}{9}% poincare-series-and-near-orthogonality
\chapters@set{20261124T223000}{10}% elliptic-curves-for-the-proof
\chapters@set{20261125T009500}{11}% fock-model-and-theta-intertwiners
\chapters@set{20261124T223100}{12}% complex-multiplication-and-inert-reduction
\chapters@set{20261125T008000}{13}% one-cm-determinant
\chapters@set{20261124T214500}{14}% supersingular-theta-divisibility
\chapters@set{20261124T214600}{15}% the-isotropic-plane
\chapters@set{20261125T013700}{16}% pairings-on-a-formal-group-of-height-two
\chapters@set{20261125T013500}{17}% hecke-nondegeneracy-on-the-tree
\chapters@set{20261126T030200}{18}% the-bound-is-attained
\chapters@set{20261125T009000}{19}% kronecker-limit-and-chowla-selberg
\chapters@set{20261125T002000}{20}% definite-theta-lifts
\chapters@set{20261125T015500}{21}% siegel-weil-in-the-smallest-case
\chapters@set{20261125T002100}{22}% definite-theta-lifts-and-norms
\chapters@set{20261125T016500}{23}% hecke-operators-and-theta-lifts
\chapters@set{20261125T016700}{24}% divisibility-along-the-chain
\chapters@set{20261125T004000}{25}% denominators-and-sha
\chapters@set{20261125T004100}{26}% tate-twists-and-sha
\chapters@set{20261126T020430}{27}% the-minus-part-by-gauss-sums
\chapters@set{20261125T014100}{28}% congruence-modules
\chapters@set{20261125T014500}{29}% eisenstein-congruences-and-selmer-classes
\chapters@set{20261125T014700}{30}% pseudocharacters-and-the-length-bound
\chapters@set{20261125T011000}{31}% galois-representations-black-box
\chapters@set{20261126T030300}{32}% the-data
\chapters@set{20261126T020210}{33}% bessel-poincare-series-and-their-gram-operators
\chapters@set{20261126T030400}{34}% the-seesaw-identity
\chapters@set{20261126T030500}{35}% divisibility-of-the-composite-lift
\chapters@set{20261126T030600}{36}% theorem-a-short-forms
\chapters@set{20261126T030700}{37}% the-determinant-of-a-chain
\chapters@set{20261126T030800}{38}% the-cost-of-one-lift
\chapters@set{20261126T030900}{39}% denominators-of-projections-to-lifts
\chapters@set{20261126T031000}{40}% twists-with-few-congruences
\chapters@set{20261126T031100}{41}% theorem-c-and-the-main-inequality
\providecommand\chapnumber[1]{%
  \@ifundefined{chapters@#1}%
    {\PackageError{chapters}{Sheet #1 is not listed in the index}%
      {List it in 20261125T012000--grh-workbook__guide.org\MessageBreak
       and run scripts/make-chapters.py.}}%
    {\edef\thechap{\@nameuse{chapters@#1}}}}
\makeatother
\chapnumber{<sheet ID>}
\makeatletter
\providecommand\chapters@set[2]{\@namedef{chapters@#1}{#2}}
\chapters@set{20261125T001000}{1}% zeros-to-ramp-sums
\chapters@set{20261125T001100}{2}% inert-primes-and-estimate-b
\chapters@set{20261125T005000}{3}% galois-fourier-and-averaging
\chapters@set{20261125T005500}{4}% columns-from-a-sphere
\chapters@set{20261125T000000}{5}% slope-method
\chapters@set{20261126T030100}{6}% a-twist-for-each-prime
\chapters@set{20261125T003000}{7}% the-chain-and-the-square-root
\chapters@set{20261125T010000}{8}% rank-lemma-and-zero-estimate
\chapters@set{20261125T010500}{9}% poincare-series-and-near-orthogonality
\chapters@set{20261124T223000}{10}% elliptic-curves-for-the-proof
\chapters@set{20261125T009500}{11}% fock-model-and-theta-intertwiners
\chapters@set{20261124T223100}{12}% complex-multiplication-and-inert-reduction
\chapters@set{20261125T008000}{13}% one-cm-determinant
\chapters@set{20261124T214500}{14}% supersingular-theta-divisibility
\chapters@set{20261124T214600}{15}% the-isotropic-plane
\chapters@set{20261125T013700}{16}% pairings-on-a-formal-group-of-height-two
\chapters@set{20261125T013500}{17}% hecke-nondegeneracy-on-the-tree
\chapters@set{20261126T030200}{18}% the-bound-is-attained
\chapters@set{20261125T009000}{19}% kronecker-limit-and-chowla-selberg
\chapters@set{20261125T002000}{20}% definite-theta-lifts
\chapters@set{20261125T015500}{21}% siegel-weil-in-the-smallest-case
\chapters@set{20261125T002100}{22}% definite-theta-lifts-and-norms
\chapters@set{20261125T016500}{23}% hecke-operators-and-theta-lifts
\chapters@set{20261125T016700}{24}% divisibility-along-the-chain
\chapters@set{20261125T004000}{25}% denominators-and-sha
\chapters@set{20261125T004100}{26}% tate-twists-and-sha
\chapters@set{20261126T020430}{27}% the-minus-part-by-gauss-sums
\chapters@set{20261125T014100}{28}% congruence-modules
\chapters@set{20261125T014500}{29}% eisenstein-congruences-and-selmer-classes
\chapters@set{20261125T014700}{30}% pseudocharacters-and-the-length-bound
\chapters@set{20261125T011000}{31}% galois-representations-black-box
\chapters@set{20261126T030300}{32}% the-data
\chapters@set{20261126T020210}{33}% bessel-poincare-series-and-their-gram-operators
\chapters@set{20261126T030400}{34}% the-seesaw-identity
\chapters@set{20261126T030500}{35}% divisibility-of-the-composite-lift
\chapters@set{20261126T030600}{36}% theorem-a-short-forms
\chapters@set{20261126T030700}{37}% the-determinant-of-a-chain
\chapters@set{20261126T030800}{38}% the-cost-of-one-lift
\chapters@set{20261126T030900}{39}% denominators-of-projections-to-lifts
\chapters@set{20261126T031000}{40}% twists-with-few-congruences
\chapters@set{20261126T031100}{41}% theorem-c-and-the-main-inequality
\providecommand\chapnumber[1]{%
  \@ifundefined{chapters@#1}%
    {\PackageError{chapters}{Sheet #1 is not listed in the index}%
      {List it in 20261125T012000--grh-workbook__guide.org\MessageBreak
       and run scripts/make-chapters.py.}}%
    {\edef\thechap{\@nameuse{chapters@#1}}}}
\makeatother
\chapnumber{20261126T051100}
\renewcommand{\thesection}{\thechap.\arabic{section}}

\theoremstyle{definition}
\newtheorem{problem}{Problem}[section]
\newtheorem{definition}{Definition}[section]
\theoremstyle{plain}
\newtheorem{theorem}{Theorem}[section]
\theoremstyle{remark}
\newtheorem*{leanremark}{In the Lean formalization}

\newcommand{\Z}{\mathbb{Z}}
\newcommand{\Q}{\mathbb{Q}}
\newcommand{\R}{\mathbb{R}}
\newcommand{\C}{\mathbb{C}}
\newcommand{\Aut}{\operatorname{Aut}}
\newcommand{\Nm}{\operatorname{Nm}}
\newcommand{\pr}{\operatorname{pr}}
\newcommand{\im}{\operatorname{im}}
\newcommand{\val}{\operatorname{val}}
\newcommand{\alg}{\mathrm{alg}}
\newcommand{\hdir}{h_{\mathrm{dir}}}
\newcommand{\sfin}{s_{\mathrm{fin}}}
\newcommand{\stot}[1]{s^{\mathrm{tot}}_{\mathrm{fin},#1}}
\newcommand{\ctest}{c_{\mathrm{test}}}
\newcommand{\catt}{c_{\mathrm{att}}}
\newcommand{\cvac}{c_{\mathrm{vac}}}
\newcommand{\cl}{\mathrm{cl}}
\newcommand{\Jet}{\mathrm{Jet}}
\newcommand{\Cl}{\mathrm{Cl}}
\newcommand{\Gm}{G_{\mathcal M}}
\newcommand{\hb}{{\mathbf{h}}}
\newcommand{\tl}[1]{\texttt{#1}}
\setlength{\emergencystretch}{3em}

\title{Test forms and the height of a Hecke ratio}
\author{}
\date{November 26, 2026}

\begin{document}
\maketitle
\refstepcounter{chap}\label{chap}

\paragraph{Prerequisites.}
Characters of finite abelian groups; \(\C_\ell\), \(\lvert\ell\rvert_\ell = 1/\ell\).  Chapter~\ref{sd-chap}, Section~\ref{sd-sec:heights}: isomorphisms \(\iota \mid \ell\), \(\lvert x\rvert_\iota := \lvert\iota(x)\rvert_\ell\), families and the height \(\mathrm{ht}\), Problems~\ref{sd-pr:height-mul} and~\ref{sd-pr:height-easy}, and the Fock norm of Problem~\ref{sd-pr:fock}.  Chapter~\ref{cg-chap}: congruence modules.  Chapter~\ref{tw-chap}, for the existence of regular families only.  Chapter~\ref{da-chap}: the data, \(\sfin\), \(s_\infty\), \(\hdir\) and the budget \(B = U/(mk)\) (Definition~\ref{da-def:budget}).  Chapter~\ref{ob-chap}: twists, \(\mathcal M\), \(\chi_M\), \(\upsilon\), \(\cvac\), and critical pairs \(P = (\nu, z)\) with their ratios \(R_P(\xi, \epsilon)\) (Definition~\ref{ob-def:pair}).  Chapter~\ref{dl-chap}: local prime sums.  Forms integral at \(\eta\) (Definition~\ref{ob-def:gauge}).  Chapter~\ref{col-chap}: two operations on heights, Problems~\ref{col-pr:of-mul} and~\ref{col-pr:height-ops}.  Chapter~\ref{l1-chap}: false lifts and clearing.  Chapter~\ref{l3-chap}: rows \(\lambda_P\) (Definition~\ref{l3-def:row}), main pairs and empty rows (Definition~\ref{l3-def:main}), \(\phi_{P,\xi}\), \(\alg_P(\xi)\) (Definition~\ref{l3-def:alg}) and Theorem~\ref{l3-thm:pair}.  Chapter~\ref{l4-chap}: \(\stot p\), the twelve copies, Theorem~\ref{l4-thm:attained}.

\medskip
This chapter closes the cost of one lift, whose statement chapter is Chapter~\ref{col-chap}.  There the height of the family \(Z_j(1)/Z_j(\theta_\iota)\), for a regular family \(\theta\) of \(m\)-fold twists, needs among other inputs a height bound for \(\iota \mapsto R_P(\vartheta_\iota)^{-1}\), \(\vartheta_\iota := \theta_\iota^j\), for the pairs \(P\) main for \(j\) in the factorization of the Hecke value product.  By Theorem~\ref{l3-thm:pair}, \(R_P(\vartheta)^{-1} = (\Nm Q_\vartheta)^{2z-1}\lVert\phi_{P,1}\rVert^2/\lVert\phi_{P,\vartheta}\rVert^2\); the untwisted squared norm has a height bound, and \(\alg_P(\vartheta)\) is integral at \(\ell \nmid \Nm Q_\vartheta\).  How does the proof bound \(\lvert\lVert\phi_{P,\vartheta_\iota}\rVert^2\rvert_\iota\) from below?  With \emph{test forms}: the lower bound~\eqref{eq:test-pairing} for their pairings with \(\phi_{P,\vartheta_\iota}\) combines the sharp jet of Chapter~\ref{l4-chap} with an auxiliary Hecke ratio.  We begin with a line in \(\C^2\).

\section{Two tests and a line}\label{sec:line}

Let \(V\) be a finite-dimensional complex inner product space, conjugate-linear in the first variable, let \(0 \ne v \in V\), and let \(\pr\) be the orthogonal projection onto \(\C v\).

\begin{problem}\label{pr:line}
For all \(\Xi, \Xi' \in V\),
\[
  \langle \pr\Xi, \pr\Xi'\rangle = \frac{\langle\Xi, v\rangle\langle v, \Xi'\rangle}{\lVert v\rVert^2} .
\]
\end{problem}

\begin{problem}\label{pr:line-toy}
In \(V = \C^2\) with the standard inner product let \(v = (1,2)\), \(\Xi = (1,0)\) and \(\Xi' = (0,1)\).  Then \(\langle\Xi,v\rangle\langle v,\Xi'\rangle = 2\) and \(\langle\pr\Xi, \pr\Xi'\rangle = 2/5\), so \(\lvert\langle\pr\Xi,\pr\Xi'\rangle\rvert_\iota = 5\) for every \(\iota \mid 5\), although \(v\), \(\Xi\) and \(\Xi'\) have integral coordinates.
\end{problem}

So projecting integral vectors to a line can create a denominator, a congruence in the sense of Chapter~\ref{cg-chap}.  Given an upper bound for \(\lvert\langle\pr\Xi,\pr\Xi'\rangle\rvert_\iota\), a lower bound for \(\lvert\langle\Xi,v\rangle\langle v,\Xi'\rangle\rvert_\iota\) gives one for \(\lvert\lVert v\rVert^2\rvert_\iota\).  For a critical pair \(P = (\nu, z)\) with \(z \ge 2\), Theorem~\ref{l3-thm:pair}(i) states that for every tame character \(\xi\)
\[
  R_P(\xi) = \frac{\lVert\phi_{P,\xi}\rVert^2/(\Nm Q_\xi)^{2z-1}}{\lVert\phi_{P,1}\rVert^2} .
\]

\begin{problem}\label{pr:five}
Let \(P = (\nu, z)\) be a critical pair with \(z \ge 2\), \(\vartheta\) a twist, \(\pr_{P,\vartheta}\) the orthogonal projection onto \(\C\phi_{P,\vartheta}\), \(\Xi, \Xi'\) forms with \(\langle\Xi,\phi_{P,\vartheta}\rangle\langle\phi_{P,\vartheta},\Xi'\rangle \ne 0\), and \(P'\) a critical pair with \(R_{P'}(\vartheta) \ne 0\).  Then
\[
  R_P(\vartheta)^{-1} = (\Nm Q_\vartheta)^{2z-1} \cdot \lVert\phi_{P,1}\rVert^2 \cdot \langle\pr_{P,\vartheta}\Xi, \pr_{P,\vartheta}\Xi'\rangle \cdot R_{P'}(\vartheta)^{-2} \cdot \frac{R_{P'}(\vartheta)^2}{\langle\Xi,\phi_{P,\vartheta}\rangle\langle\phi_{P,\vartheta},\Xi'\rangle} .
\]
\end{problem}

In the application \(P\) is main for \(j\), \(\vartheta = \theta_\iota^j\) for a regular family \(\theta\), \(\Xi\) and \(\Xi'\) are test forms and \(P'\) is ramified; Problem~\ref{pr:sum} bounds the five factors.

\begin{definition}\label{def:test}
Let \(P = (\nu, z)\) be a critical pair with \(z \ge 2\), \(\vartheta\) a twist and \(\ell\) a prime.  We define a \emph{test form for \(\ell\)} (for \(P\) and \(\vartheta\)) to be a form \(\Xi\) in the space of forms of \(\phi_{P,\vartheta}\), on the definite space of rank \(2z\) (Definition~\ref{l3-def:char-lift}), such that:
\begin{enumerate}
\item \(\Xi\) is integral at every ring homomorphism \(\eta\colon \C \to \C_\ell\) (Definition~\ref{ob-def:gauge});
\item \(\Xi\) is orthogonal to the false lifts of the rank-one lift: for every finite place \(w\) of \(E\) with \(\bar w \ne w\) that is not a place of the lift or the conjugate of one, \(\Xi\) is orthogonal to every form in the image of \(e_{X-1}\) that is killed by the Hecke polynomial at \(w\) at its test value (as in Definition~\ref{l1-def:false}) and by the descent, the adjoint of the lift.  Here \(e_{X-1}\) is the idempotent with label \(X - 1\) at every chosen place; since the source has rank one, it is the mean over \(\prod_{w \in \mathcal M} k_w^\times\).
\end{enumerate}
\end{definition}

Let \(\theta = (\theta_\iota)\) be a family of \(m\)-fold twists, one for every prime \(\ell\) and every \(\iota \mid \ell\).

\begin{definition}\label{def:small-ratios}
The family \(\theta\) has \emph{small ratios} if (i) \(\ell\) is prime to \(\Nm Q_{\theta_\iota}\) for every \(\iota \mid \ell\), and (ii) for every \(1 \le j \le m\), every ramified critical pair \(P\) and \(\epsilon = \pm 1\), the family \(\iota \mapsto R_P(\theta_\iota^j, \epsilon)^{-1}\) has height at most \(x_1\).
\end{definition}

\begin{definition}\label{def:regular}
The family \(\theta\) is \emph{regular} if it has small ratios and (iii) for every \(1 \le j \le m\) there are reals \(\cl(\ell)\), one for each prime \(\ell\), with \(\sum_{\ell \in P_0} \cl(\ell) \le B\) for every finite set \(P_0\) of primes, such that for every \(\iota \mid \ell\), with \(\vartheta := \theta_\iota^j\): \(\cl(\ell)\) clears the projection at \((\iota, \vartheta, j)\) (Definition~\ref{l1-def:clears}); and for every critical pair \(P\) main for \(j\) and any two test forms \(\Xi\), \(\Xi'\) for \(\ell\) for \(P\) and \(\vartheta\), \(\lvert\langle\pr_{P,\vartheta}\Xi, \pr_{P,\vartheta}\Xi'\rangle\rvert_\iota \le e^{\cl(\ell)}\).
\end{definition}

That regular families exist is Theorem~\ref{tw-thm:regular}.

\section{A large pairing from characters}\label{sec:large}

Where does a test form with a large pairing come from?  In the proof of~\eqref{eq:test-pairing} its pairing with \(\phi_{P,\vartheta}\) is a scalar multiple of a normalized character sum of a theta kernel with twelve coordinates at each \(\tau\); after scaling, the value of the kernel at the identity is the coefficient estimated in Theorem~\ref{l4-thm:attained}.  Two elementary facts carry the passage.

\begin{problem}\label{pr:orthogonality}
Let \(\Cl\) be a finite abelian group of order \(N_c\), \(\mathrm{kn}\colon \Cl \to \C\), and for a character \(\psi\) of \(\Cl\) let \(\mathrm{av}_\psi := N_c^{-1}\sum_{g \in \Cl}\psi(g)\,\mathrm{kn}(g)\).  Then \(\sum_\psi \mathrm{av}_\psi = \mathrm{kn}(1)\), the sum over all characters, and for every prime \(\ell\) and every ring homomorphism \(\eta\colon\C\to\C_\ell\) some \(\psi\) has \(\lvert\eta(\mathrm{av}_\psi)\rvert \ge \lvert\eta(\mathrm{kn}(1))\rvert\).  If \(N_c = 2\) and \(\mathrm{kn}\) is the indicator of \(1\), both \(\mathrm{av}_\psi\) equal \(1/2\), of \(2\)-adic absolute value \(2 \ge \lvert\mathrm{kn}(1)\rvert_2\).
\end{problem}

\begin{problem}\label{pr:pigeonhole}
Let \(G\) be a finite group, \(\mathcal A_3, \mathcal A_1\colon G \to \C\), and \(L_* \in \C\) with \(L_*^{\#G} = \prod_{\chi \in G} \mathcal A_3(\chi)\mathcal A_1(\chi)\).  Let \(\eta\colon\C\to F_0\) be a ring homomorphism into a field with an absolute value.  Then for every \(\psi \in G\) there are \(\chi_3, \chi_1 \in G\) with \(\chi_3\chi_1 = \psi\) and \(\lvert\eta(\mathcal A_3(\chi_3)\mathcal A_1(\chi_1))\rvert \ge \lvert\eta(L_*)\rvert\).  For \(G = \{1, \psi\}\), \(\mathcal A_3 = (1, 9)\), \(\mathcal A_1 = (9, 1)\) and \(L_* = 9\), exactly one factorization of \(\psi\) satisfies the inequality for \(\lvert\cdot\rvert_3\).
\end{problem}

For a main pair the proof of~\eqref{eq:test-pairing} splits the \(2z\) coordinates into blocks of sizes \(12\), \(4\), \(4\) and \(2z - 20\): the jet lives on the first block, the twist on the second, and the last two carry auxiliary lifts.

\begin{problem}\label{pr:blocks}
For an integer \(z \ge 2\), \(12 + 4 + 4 + \max(2z - 20, 0) = 2z\) if and only if \(z \ge 10\).  For a pair main for \(j\), \(10000 \le m < z \le r\) and \(4 \le 2z - 20 \le t\).
\end{problem}

Let \(P\) be main for \(j\), \(\lambda := \lambda_P\) its row, and \(\hb = (h_{\tau,i})\) a multi-index on the twelve copies with \(\sum_i h_{\tau,i} = \lambda(\tau)\) for every \(\tau\) (Chapter~\ref{l4-chap}).  For the twelve coordinates of the first block the proof uses an integer vector in the Weyl representation of the lift; over every field of characteristic zero it is fixed by every family of invertible matrices fixing those twelve basis vectors at each \(\tau\), and its Fock polynomial in these coordinates is \(T_\hb X^\hb\), with
\[
  T_\hb := \prod_{\tau\in\Phi}\frac{\lambda(\tau)!}{\prod_{i=1}^{12} h_{\tau,i}!} .
\]
This is read from the proof; not proved here.  Give the polynomials in the \(X_{\tau,i}\) the Fock inner product of Problem~\ref{sd-pr:fock} in each variable: \(\lVert X^\hb\rVert^2 = \prod_{\tau,i}h_{\tau,i}!\), and monomials are orthogonal.

\begin{problem}\label{pr:multinomial}
\(T_\hb\) is the coefficient of \(X^\hb\) in \(\prod_\tau (X_{\tau,1} + \dots + X_{\tau,12})^{\lambda(\tau)}\), and \(T_\hb \in \Z\).  For a polynomial \(f = \sum_{\hb'} a_{\hb'} X^{\hb'}\), \(\langle T_\hb X^\hb, f\rangle = \prod_\tau \lambda(\tau)! \cdot a_\hb\).  For one place, \(\lambda = 13\) and \(\hb = (2, 1, \dots, 1)\), \(T_\hb = 13!/2\).
\end{problem}

Theorem~\ref{l4-thm:attained}, for \(\lambda\) with \(\kappa(\tau) \le \lambda(\tau) \le \kappa(\tau) + x^{42}\) for every \(\tau\), gives \(\catt \in \C^\times\) with \(\lvert\sigma(\catt)\rvert \le \exp(fmb/k)\) for every \(\sigma \in \Aut(\C)\) such that for the modulus \(\mathfrak m := (\pi_M)\) (Definition~\ref{c3-def:modulus}), every prime \(p\) and every ring homomorphism \(\eta\colon\C\to\C_p\) there are \(\hb\) with \(\sum_i h_{\tau,i} = \lambda(\tau)\) for every \(\tau\) and \(y_0 \in (\mathcal O_E/\mathfrak m)^{12}\) with
\[
  \val\Bigl(\eta\Bigl(\frac{\cvac^{12}}{\catt}\prod_{\tau\in\Phi}\lambda(\tau)!\cdot[z^\hb]\Jet_{y_0}\Bigr)\Bigr) \le \stot p .
\]
For a pair \(P\) main for \(j\), the row \(\lambda_P\) satisfies these bounds (Problem~\ref{l3-pr:row-range}).

Now the objects read from the proof.  In this sheet \(\Gm\) is the finite group of pairs \(\psi\) consisting of a continuous unitary character of the norm-one ideles of \(E\), trivial on the norm-one elements of \(E^\times\), of weight \(0\), and its tame part, which is trivial off \(\mathcal M\); \(\psi(g)\) is the value of the character at the norm-one idele given by a finite adelic point \(g\) of \(U(1)\).  For every \(\psi \in \Gm\) and every even \(n\) with \(4 \le n \le t\) there is a critical pair \((\nu, n/2)\) with empty row, of \emph{rank} \(n\) (the rank \(2z\) of the space of its lift, Definition~\ref{l3-def:char-lift}), with the tame part of \(\psi\), whose character becomes that of \(\psi\) after multiplication by \(\upsilon^{n/2}\).  For a pair \(P = (\nu_P, z)\) main for \(j\), let \(P^{(4)}_\psi\) and \(P^{(2z-20)}_\psi\) be such pairs of ranks \(4\) and \(2z - 20\), and fix \(L_*\) with
\[
  L_*^{\#\Gm} = \prod_{\psi \in \Gm}\alg_{P^{(4)}_\psi}(1)\,\alg_{P^{(2z-20)}_\psi}(1) .
\]
For a pair \(P_0\) with empty row, \(\alg_{P_0}(1)\) is a positive real number (Chapter~\ref{l3-chap}), and
\[
  \lvert\sigma(\alg_{P_0}(1))\rvert \le e^{x^{34}} \qquad\text{for every } \sigma \in \Aut(\C) ;
\]
this bound is read from the proof; not proved here.  Hence \(L_* \ne 0\) and \(\lvert\sigma(L_*)\rvert \le e^{2x^{34}}\) for every \(\sigma\in\Aut(\C)\).  The classes of the finite adelic points of \(U(1)\) modulo the rational points and the principal level at the chosen places form a finite abelian group \(\Cl\), of order \(N_c\); every character of \(\Cl\) has the form \(g \mapsto \psi(g)\) for some \(\psi \in \Gm\).  Fix representatives \(g_c\) of the classes \(c \in \Cl\).

The exact pairing.  Let \(P\) be main for \(j\), \(P'\) a critical pair with empty row of rank four, \(\ell\) a prime, \(\eta\colon\C\to\C_\ell\) a ring homomorphism and \(\vartheta\) a twist with \(\ell \nmid \Nm Q_\vartheta\).  Let \(\hb\) and \(y_0\) be given by Theorem~\ref{l4-thm:attained} for \(\lambda = \lambda_P\), \((\ell, \eta)\) and the modulus \((\pi_M)\).  Write \(\nu_{P'}\) for the character of \(P'\) and put \(\chi_0 := (\nu_P/\nu_{P'})\,\upsilon^{z-2}\).  Let \(\mathrm{kn}(g)\) be the scalar theta kernel on the twelve-block, with characters \(1\) and \(\chi_0\), polynomial \(X^\hb\) and the indicator of \(y_0\), evaluated at the point of \(U(1)\) with finite part \(g\) and trivial archimedean part, and at the identity of \(U(12)\).  There is a number \(c_{\mathrm{kn}} \ne 0\) such that \(c_{\mathrm{kn}}\,\mathrm{kn}(1) = \cvac^{12}\catt^{-1}\prod_\tau\lambda_P(\tau)!\cdot[z^\hb]\Jet_{y_0}\), so that \(\lvert\eta(c_{\mathrm{kn}}\,\mathrm{kn}(1))\rvert \ge \ell^{-\stot\ell}\); and for every \(\psi \in \Gm\), every factorization \(\psi = \psi_3\psi_1\) in \(\Gm\) and every finite adelic point \(g_0\) of \(U(1)\) there is a test form \(\Xi\) for \(\ell\) for \(P\) and \(\vartheta\) with
\[
  N_c^2\,\langle\Xi, \phi_{P,\vartheta}\rangle = \catt\,\alg_{P'}(1)\,\alg_{\psi_3,\psi_1}\,R_{P'}(\vartheta)\,c_{\mathrm{kn}}\,\mathrm{av}_{\psi,g_0},
  \qquad
  \mathrm{av}_{\psi,g_0} := \frac{1}{N_c}\sum_{c \in \Cl}\psi(g_cg_0)\,\mathrm{kn}(g_cg_0) ,
\]
where \(\alg_{\psi_3,\psi_1} := \alg_{P^{(4)}_{\psi_3}}(1)\,\alg_{P^{(2z-20)}_{\psi_1}}(1)\).  Here \(\Xi\) is a class sum, against the lift of \(P'\) at \(\vartheta\) and the untwisted lifts of the two auxiliary pairs, of a point form with the integer vector above.  \emph{Status.}  We take this identity as an input.  The Lean proof obtains it for its chosen character, factorization and translation.  The stated form follows from the pointwise kernel identity by summing over translated classes, as in the final remark.

\begin{problem}\label{pr:test-pairing}
Let \(P\) be main for \(j\) and \(P'\) a critical pair with empty row of rank four, and put \(\ctest := \catt\,\alg_{P'}(1)\,L_*\).  Then \(\ctest \ne 0\), \(\lvert\sigma(\ctest)\rvert \le \exp(fmb/k + 3x^{34})\) for every \(\sigma\in\Aut(\C)\), and for every prime \(\ell\), every ring homomorphism \(\eta\colon\C\to\C_\ell\) and every twist \(\vartheta\) with \(\ell\nmid\Nm Q_\vartheta\) there is a test form \(\Xi\) for \(\ell\) for \(P\) and \(\vartheta\) with
\begin{equation}\label{eq:test-pairing}
  \lvert\eta(\ctest\,R_{P'}(\vartheta))\rvert\cdot\ell^{-\stot\ell} \le \lvert\eta\langle\Xi, \phi_{P,\vartheta}\rangle\rvert .
\end{equation}
\end{problem}

\section{The height of a Hecke ratio}\label{sec:height}

Let \(\theta\) be regular, \(1 \le j \le m\), \(P = (\nu, z)\) a critical pair main for \(j\), and put \(\vartheta_\iota := \theta_\iota^j\).  Fix a ramified critical pair \(P'\) with empty row of rank four; the proof takes one whose tame part is that of \(\chi_M\) and whose character becomes \(\chi_M\) after multiplication by \(\upsilon^2\).  Let \(\ctest\) be as in Problem~\ref{pr:test-pairing} for \(P\) and \(P'\).  The ratio \(R_{P'}(\vartheta_\iota)\) is a positive real number.

\begin{problem}\label{pr:two-tests}
Let \(\iota \mid \ell\).  There are test forms \(\Xi_\iota\), \(\Xi'_\iota\) for \(\ell\) for \(P\) and \(\vartheta_\iota\) such that, with \(y_\iota := \langle\Xi_\iota, \phi_{P,\vartheta_\iota}\rangle\langle\phi_{P,\vartheta_\iota}, \Xi'_\iota\rangle\),
\[
  \lvert\ctest\,\bar\ctest\,R_{P'}(\vartheta_\iota)^2\rvert_\iota\cdot\ell^{-2\stot\ell} \le \lvert y_\iota\rvert_\iota ;
\]
in particular \(y_\iota \ne 0\).
\end{problem}

For each \(\iota\) fix such \(\Xi_\iota\), \(\Xi'_\iota\), and let \(y_\iota\) be the resulting product.

\begin{problem}\label{pr:prime-sums}
For every finite set \(P_0\) of primes,
\[
  \sum_{p \in P_0} 2\,\stot p\log p \le \sum_{\tau\in\Phi}\sfin(\kappa(\tau), \tau) .
\]
\end{problem}

\begin{problem}\label{pr:sizes}
For the data, \(fmb/k + 3x^{34} \le B\) and \(x_1 \le B\).
\end{problem}

\begin{problem}\label{pr:pairing-height}
The family \(\iota \mapsto R_{P'}(\vartheta_\iota)^2/y_\iota\) has height at most \(2B + \sum_{\tau\in\Phi}\sfin(\kappa(\tau),\tau)\).
\end{problem}

\begin{problem}\label{pr:sum}
In Problem~\ref{pr:five} at \(\vartheta_\iota\), with \(\Xi_\iota, \Xi'_\iota\) and \(P'\), the five factors have \(\lvert(\Nm Q_{\vartheta_\iota})^{2z-1}\rvert_\iota = 1\) and, as families in \(\iota\), heights at most \(B - \sum_\tau s_\infty(\kappa(\tau))\), \(B\), \(2x_1\) and \(2B + \sum_\tau \sfin(\kappa(\tau),\tau)\) for the last four.  The sum of these four bounds is \(4B + 2x_1 - \hdir \le 6B - \hdir\).
\end{problem}

Problems~\ref{pr:five} and~\ref{pr:two-tests}--\ref{pr:sum} prove the following.

\begin{theorem}\label{thm:main-pair}
Let \(\theta\) be regular, \(1 \le j \le m\), and \(P\) a critical pair main for \(j\).  Then the family \(\iota \mapsto R_P(\theta_\iota^j)^{-1}\) has height at most \(6B - \hdir\).
\end{theorem}

\begin{problem}\label{pr:negative}
For the data, \(6B < \hdir\).  So the bound of Theorem~\ref{thm:main-pair} is negative.
\end{problem}

We have proved Theorem~\ref{thm:main-pair} from Theorems~\ref{l3-thm:pair} and~\ref{l4-thm:attained} and the exact pairing of Section~\ref{sec:large}, taken as an input with the auxiliary pairs, the conjugate bound for their algebraic parts, the integer vector on the twelve-block, and the finite class group \(\Cl\) with the lifting of its characters to \(\Gm\); the bounds on ramified ratios and on denominators of projections enter as the definition of regular families.  By Problem~\ref{l3-pr:hecke-product}, at \(\vartheta_\iota\) the inverse ratio of Hecke value products is, up to a unit at \(\ell\), the square of a product of \(k/2\) factors \(R_P(\vartheta_\iota)^{-1}\), \(P\) main, and \(5m\) factors \(R_{P'}(\vartheta_\iota, \epsilon)^{-1}\), \(P'\) ramified, with signs depending only on \(j\); so its height is at most \(k(6B - \hdir) + 10mx_1\) (Problem~\ref{col-pr:hecke-height}).

\begin{leanremark}
\raggedright In the Lean proof \cite{GL} (\tl{grh/src/New/}), Theorem~\ref{thm:main-pair} is \tl{Data.heightLe\_inv\_ratio}, whose proof also chooses \(P'\); \eqref{eq:test-pairing} is \tl{Data.exists\_isTestForm\_le\_norm\_inner\_trivialLift}, in whose proof are the group \(\Gm\), the blocks, the classes \(\Cl\) and the value \(c_{\mathrm{kn}}\,\mathrm{kn}(1)\) (step \tl{h\_lower}); its step \tl{h\_inner} is the exact pairing for the choices made there.  The exact pairing in general follows from \tl{Data.exists\_isTestForm\_inner\_kernelForm\_eq}, \tl{Data.modKernel\_div\_apply} and \tl{LValueData.trivialLift\_eq\_stdClassSum\_translate}.  Test forms are \tl{DescentPair.IsTestForm}; Definitions~\ref{def:small-ratios} and~\ref{def:regular} are \tl{Data.HasSmallRatios} and \tl{Data.IsRegularFamily}.  The auxiliary pairs come from \tl{Data.exists\_hasEmptyRow}, the bound \(e^{x^{34}}\) is \tl{Data.CriticalPair.norm\_map\_integralAlgebraicPart\_le}, \(L_*\) is from \tl{Data.exists\_root\_prod\_integralAlgebraicPart}, the integer vector is \tl{Rows.exists\_int\_weylRep\_invariant\_fockPoly\_eq\_smul\_monomial}, and the lifting of characters of \(\Cl\) is in \tl{exists\_norm\_le\_norm\_twisted\_classSum}.
\end{leanremark}

\section*{Hints}

\subsection*{\S\ref{sec:line}}
\begin{description}
\item[\ref{pr:line}] \(\pr\Xi = (\langle v,\Xi\rangle/\lVert v\rVert^2)\,v\), and the first slot is conjugate-linear.
\item[\ref{pr:five}] Theorem~\ref{l3-thm:pair}(i) at \(\xi = \vartheta\), and Problem~\ref{pr:line} for \(\lVert\phi_{P,\vartheta}\rVert^2\).
\end{description}

\subsection*{\S\ref{sec:large}}
\begin{description}
\item[\ref{pr:orthogonality}] \(\sum_\psi\psi(g)\) is \(N_c\) or \(0\).  For the second statement, \(\C_\ell\) is ultrametric.
\item[\ref{pr:pigeonhole}] For fixed \(\psi\), compare the product of \(\lvert\eta(\mathcal A_3(\psi\chi^{-1})\mathcal A_1(\chi))\rvert\) over \(\chi\) with its largest factor.
\item[\ref{pr:multinomial}] The multinomial theorem, and Problem~\ref{sd-pr:fock} variable by variable.
\item[\ref{pr:test-pairing}] Translate the representative of the identity class to \(1\).  Which character and which factorization give the two lower bounds?  Use \(N_c \in \Z\) to control its local absolute value.
\end{description}

\subsection*{\S\ref{sec:height}}
\begin{description}
\item[\ref{pr:two-tests}] Keep the twist fixed.  Which two maps \(\C \to \C_\ell\) account for the two orders of the arguments of the inner product?
\item[\ref{pr:prime-sums}] Problems~\ref{dl-pr:part-at-p} and~\ref{da-pr:sfin}; all terms are nonnegative.
\item[\ref{pr:sizes}] \(B = fxb\log^3 b/(mk)\).  Use Definition~\ref{da-def:data} and Problem~\ref{da-pr:m-big}: \(m^2 \le 4x\), \(k \le r \le 3x^5\), \(x^{10^6} \le x_1\) and \(x_1^{10^6} \le b\).
\item[\ref{pr:pairing-height}] Problem~\ref{col-pr:of-mul}: which multiplier works at both embeddings?  Then Problems~\ref{pr:two-tests}, \ref{pr:prime-sums} and~\ref{pr:sizes}.
\item[\ref{pr:sum}] Match the four factors with Theorem~\ref{l3-thm:pair}(iii), clause (iii) of Definition~\ref{def:regular} with Problem~\ref{sd-pr:height-easy}(a), clause (ii) of Definition~\ref{def:small-ratios} for \(P'\) squared (Problem~\ref{sd-pr:height-mul}) and Problem~\ref{pr:pairing-height}; then Problem~\ref{col-pr:height-ops}.
\item[\ref{pr:negative}] \(\hdir \ge b\,\#\Phi_a\) (Problem~\ref{da-pr:hdir}), \(f = \#\Phi_a\,\varphi(q_0)\) (Problem~\ref{da-pr:degree}), and compare \(6\varphi(q_0)x\log^3 b\) with \(mk\).
\end{description}

\begin{thebibliography}{GL}
\bibitem[GL]{GL} L.~Sernau, \emph{GRH for Dirichlet \(L\)-functions, in Lean~4}, a shorter rewriting of the Lean proof, text by Claude, \url{https://github.com/lsernau/grh-lean}, commit \texttt{72fead1}, readable edition \texttt{grh/}.
\end{thebibliography}

\end{document}
